\chapter{Matrix Lie Groups}

\begin{center}\itshape
A full turn ends at the identity rotation. The path that makes the turn
still remembers something: it cannot be shrunk to a constant path.
Two full turns can.
\end{center}

A rotation is a matrix, and nearby rotations have nearby entries.
The rotation group is therefore both a group and a geometric space.
Its multiplication describes how motions combine; its tangent vectors
describe their initial velocities. The link between these two descriptions
is the matrix exponential.

The distinction between a rotation and a path of rotations also appears
in spin. Uhlenbeck and Goudsmit proposed electron spin in 1925, and Pauli's
1927 description uses two-component complex state vectors.
The group acting on these vectors is $SU(2)$, which covers $SO(3)$ twice.
A lifted full-turn path ends at $-I$; two turns end at $I$.
This is not a claim that the final spatial rotation differs from the
identity. Nor does an overall sign by itself distinguish physical states.
A relative sign between two coherent parts of a beam can be detected by
interference. The neutron experiments of Rauch and collaborators and of
Werner and collaborators in 1975 observed the corresponding phase shift.

We can prove the double-cover statement using the tools already developed.
First, a closed matrix subgroup turns out to be a manifold without any
smoothness assumption. Its tangent space at the identity is a vector
space closed under a bracket. Then homomorphisms become linear maps on
these tangent spaces. For $SU(2)$ and $SO(3)$, the tangent spaces agree
while the groups have different global topology.

\section{Matrix groups and their Lie algebras}

Matrix sizes are positive integers unless the zero-size convention is
explicitly mentioned. Complex matrices may be read as real matrices of
twice the size: multiplication by $a+ib$ acts on $\R^2$ as
$\left(\begin{smallmatrix}a&-b\\b&a\end{smallmatrix}\right)$.
All tangent spaces and dimensions below are real.
Use the operator norm $\|X\|=\sup_{|v|=1}|Xv|$, which satisfies
$\|XY\|\le\|X\|\|Y\|$.

\begin{definition}
A matrix Lie group is a subgroup $G$ of $GL_n(\R)$ or $GL_n(\C)$
that is closed in $GL_n$ with its usual topology.
Its Lie algebra is the set
\[
 \mathfrak g=\{X\in M_n:e^{tX}\in G\text{ for every }t\in\R\}.
\]
\end{definition}

Closedness is relative to $GL_n$, not to all matrices: a sequence may
approach a singular matrix without violating it.
Examples include the determinant-one group $SL_n$, the orthogonal groups
$O(n)$ and $SO(n)$, the unitary groups $U(n)$ and $SU(n)$, and finite
matrix groups. Their defining equations are continuous.
The Heisenberg group of upper triangular $3\times3$ matrices with diagonal
entries one is another example.
By contrast, for irrational $\alpha$ the subgroup
$\{(e^{2\pi it},e^{2\pi i\alpha t}):t\in\R\}$ of $U(1)^2$ is dense
and not closed, as the torus flow of Figure~\ref{fig:torusflow} illustrates.
It does not meet this definition.

Volume~8, Proposition~C.13, established the matrix exponential and the
identities
\[
 e^{X+Y}=e^Xe^Y\quad(XY=YX),\qquad
 \det e^X=e^{\operatorname{tr}X},\qquad
 \frac{\dd}{\dd t}e^{tX}=Xe^{tX}.
\]
It also gives $e^{PXP^{-1}}=Pe^XP^{-1}$.
Near $I$, the exponential can be undone by the logarithmic series.

\begin{lemma}\label{lem:log}
Define $\log A=\sum_{k\ge1}(-1)^{k+1}(A-I)^k/k$ for $\|A-I\|<1$.
\begin{enumerate}
\item[(a)] In that range, $e^{\log A}=A$.
If $\|X\|<\log2$, then $\|e^X-I\|<1$ and $\log(e^X)=X$.
\item[(b)] A constant $c$ satisfies
$\|\log A-(A-I)\|\le c\|A-I\|^2$ whenever $\|A-I\|\le1/2$.
\item[(c)] The exponential is a diffeomorphism from a neighbourhood of $0$
in $M_n$ onto a neighbourhood of $I$. Its differential at $0$ is the identity.
\end{enumerate}
\end{lemma}
\begin{proof}
(a) For scalars, the corresponding exponential and logarithmic power
series compose to the identity near zero. Comparing their convergent
series gives identities of coefficients. The same coefficients apply
to powers of a single matrix; absolute convergence in the operator norm
justifies the compositions and rearrangements.
In particular,
$\|e^X-I\|\le e^{\|X\|}-1<1$ for $\|X\|<\log2$.
The scalar identity $\log e^z=z$ on this neighbourhood therefore gives
$\log(e^X)=X$ as well.
(b) The tail is bounded by
\[
 \sum_{k\ge2}\frac{\|A-I\|^k}{k}
 \le\|A-I\|^2\sum_{j\ge0}2^{-j}=2\|A-I\|^2.
\]
Thus $c=2$ works.
(c) The exponential series is smooth and
$e^X=I+X+O(\|X\|^2)$. Its differential at zero is the identity.
The inverse function theorem, Theorem~\ref{thm:inverse}, gives the local
diffeomorphism; part (a) identifies its inverse with this logarithm.
\end{proof}

A one-parameter subgroup is a homomorphism from the additive group $\R$.
Even continuity, with no derivative assumed, forces it to be exponential.

\begin{proposition}\label{prop:oneparameter}
Every continuous homomorphism $\gamma:\R\to GL_n$ has the form
$\gamma(t)=e^{tX}$ for a unique matrix $X$.
It is consequently smooth, with $X=\gamma'(0)$.
\end{proposition}
\begin{proof}
Choose $\varepsilon>0$ so that for $|t|\le\varepsilon$,
$\|\gamma(t)-I\|<1/2$ and $\|\log\gamma(t)\|<(\log2)/2$.
Set $X=(\log\gamma(\varepsilon))/\varepsilon$.
Let $B=\log\gamma(\varepsilon/2)$.
The homomorphism property gives
$e^{2B}=\gamma(\varepsilon)=e^{\varepsilon X}$.
Both exponents lie in the logarithm's inverse range, so $2B=\varepsilon X$.
Repeat the same argument at successive halves of $\varepsilon$.
It gives $\gamma(\varepsilon/2^k)=e^{\varepsilon X/2^k}$ for every $k$.
Taking positive and negative integer powers gives equality at every
$m\varepsilon/2^k$, $m\in\Z$. These points are dense in $\R$.
Continuity on both sides proves equality everywhere.
Differentiating the resulting exponential formula at zero gives
$X=\gamma'(0)$, proving uniqueness.
\end{proof}

For noncommuting matrices, the exponential of a sum is not generally the
product of the exponentials. Small products still recover both the sum
and the commutator, and closedness lets their limits remain in the group.

\begin{proposition}\label{prop:liealgebra}
For matrices $X,Y$, put $[X,Y]=XY-YX$. Then
\[
 e^{X+Y}=\lim_{m\to\infty}(e^{X/m}e^{Y/m})^m,
\]
\[
 e^{[X,Y]}=\lim_{m\to\infty}
 (e^{X/m}e^{Y/m}e^{-X/m}e^{-Y/m})^{m^2}.
\]
Consequently $\mathfrak g$ is a real vector space closed under the
commutator, and $gXg^{-1}\in\mathfrak g$ for $g\in G$, $X\in\mathfrak g$.
\end{proposition}
\begin{proof}
The first product is
$I+(X+Y)/m+O(m^{-2})$. Its logarithm is therefore
$Z_m=(X+Y)/m+O(m^{-2})$, by Lemma~\ref{lem:log}.
Its $m$th power is $e^{mZ_m}$, tending to $e^{X+Y}$.
For the second product, expand each exponential through degree two in
$s=1/m$. The first three factors give
\[
 e^{sX}e^{sY}e^{-sX}
 =I+sY+s^2([X,Y]+Y^2/2)+O(s^3).
\]
Multiplying by $e^{-sY}=I-sY+s^2Y^2/2+O(s^3)$ cancels the linear
term and both remaining $Y^2$ contributions. The product is therefore
$I+s^2[X,Y]+O(s^3)$. Its logarithm has the same leading term.
Multiplying that logarithm by $m^2$ gives $[X,Y]+O(m^{-1})$, proving
the second limit by continuity of the exponential.

If $X,Y\in\mathfrak g$, scalar multiples stay in $\mathfrak g$ by its
definition. Replace $X,Y$ by $tX,tY$ in the limits.
Every approximating product lies in $G$, and the limits are invertible,
so relative closedness in $GL_n$ gives
$e^{t(X+Y)}\in G$ and $e^{t^2[X,Y]}\in G$.
Interchanging $X,Y$ supplies the negative multiples of the bracket.
Thus every real multiple occurs, proving bracket closure.
Finally $e^{tgXg^{-1}}=ge^{tX}g^{-1}\in G$ proves conjugation invariance.
\end{proof}

A Lie algebra is a real vector space with a bilinear antisymmetric bracket
satisfying the Jacobi identity
$[X,[Y,Z]]+[Y,[Z,X]]+[Z,[X,Y]]=0$.
For matrices, expanding gives each of the six ordered triple products
twice with opposite signs. Thus the commutator has this identity.

Differentiate a group's defining equation along $e^{tX}$ at $t=0$ to
find necessary conditions on $X$. The exponential identities give the
converse conditions:
\[
 \mathfrak{sl}_n=\{X:\operatorname{tr}X=0\},\qquad
 \mathfrak{so}(n)=\{X:X^{\mathsf T}=-X\},
\]
\[
 \mathfrak u(n)=\{X:X^*=-X\},\qquad
 \mathfrak{su}(n)=\{X:X^*=-X,\ \operatorname{tr}X=0\}.
\]
For example, $(e^{tX})^*e^{tX}=I$ differentiates to $X^*+X=0$.
Conversely, if $X^*=-X$, the commuting exponentials $e^{tX^*}$ and
$e^{tX}$ multiply to $I$.
For determinant one, differentiating $e^{t\operatorname{tr}X}=1$
at $t=0$ gives $\operatorname{tr}X=0$, and that condition also suffices.
Counting free real entries gives dimensions $3$ for $\mathfrak{so}(3)$,
$3$ for $\mathfrak{su}(2)$, $n^2$ for $\mathfrak u(n)$ and $3$ for
$\mathfrak{sl}_2(\R)$.
The next theorem shows that these really are the dimensions of the groups.

\begin{theorem}[von Neumann, Cartan]\label{thm:closedsubgroup}
Let $G\subset GL_n$ be a matrix Lie group with Lie algebra $\mathfrak g$.
There are neighbourhoods $U$ of $0$ in $M_n$ and $V$ of $I$ such that
$\exp:U\cap\mathfrak g\to V\cap G$ is a homeomorphism.
The group $G$ is a smooth submanifold of $M_n$, of dimension
$\dim\mathfrak g$, with $T_IG=\mathfrak g$. Multiplication and inversion
on $G$ are smooth.
\end{theorem}
\begin{proof}
Choose a vector-space complement $W$ to $\mathfrak g$ in $M_n$.
First show that for some $r>0$,
\[
 Y\in W,\quad0<\|Y\|<r\quad\Longrightarrow\quad e^Y\notin G.
\]
If no such $r$ existed, there would be nonzero $Y_m\in W$ tending to
zero with $e^{Y_m}\in G$. Passing to a subsequence on the compact unit
sphere gives $Y_m/\|Y_m\|\to Y\in W$ with $\|Y\|=1$.
For each fixed $t\in\R$, choose integers $k_m$ with
$k_m\|Y_m\|\to t$. Then
$e^{k_mY_m}=(e^{Y_m})^{k_m}\in G$ and $k_mY_m\to tY$.
Closedness gives $e^{tY}\in G$ for every $t$, hence
$Y\in\mathfrak g\cap W=\{0\}$, a contradiction.
If $W=\{0\}$ the assertion is vacuous and any $r$ works.

Consider $\Phi(X,Y)=e^Xe^Y$ on $\mathfrak g\times W$.
Its differential at $(0,0)$ is $(X,Y)\mapsto X+Y$, an isomorphism.
The inverse function theorem gives small product neighbourhoods
$U_{\mathfrak g}\times U_W$ on which $\Phi$ is a diffeomorphism onto
an open set $V$ about $I$. Choose $U_W$ inside the radius-$r$ ball,
and $U_{\mathfrak g}$ inside a neighbourhood where $\exp$ is injective.
For $A=e^Xe^Y\in V$, if $A\in G$ then
$e^Y=e^{-X}A\in G$. Our first assertion forces $Y=0$.
Conversely $Y=0$ gives $A=e^X\in G$.
Thus in these coordinates
\[
 G\cap V=\Phi(U_{\mathfrak g}\times\{0\})=\exp(U_{\mathfrak g}).
\]
The coordinate slice is a submanifold, and its tangent space at $I$ is
$\mathfrak g$ because $d\exp_0$ is the identity.
Set $U=U_{\mathfrak g}+U_W$, an open neighbourhood in $M_n$.
Then $U\cap\mathfrak g=U_{\mathfrak g}$, giving the required
homeomorphism, in fact a diffeomorphism of these submanifolds.

Left multiplication by any $g\in G$ is an invertible linear map on
$M_n$ and carries this local description to a neighbourhood of $g$.
It follows that all of $G$ is a submanifold of the claimed dimension.
Matrix multiplication and inversion are smooth on $GL_n$;
their restrictions to this submanifold give smooth group operations.
The argument also includes $\mathfrak g=0$, when the local slice is a point.
\end{proof}

\begin{beauty}
No parametrisation of the closed subgroup was assumed.
The proof found one by separating genuine exponential directions from a
transverse complement. A hypothetical sequence of small transverse group
elements would, by taking large integer powers, create another genuine
direction. That is the contradiction that makes the chart possible.
Hilbert's fifth problem asks a broader question about topological groups.
The theorem of Gleason, Montgomery and Zippin that a Hausdorff,
second-countable locally Euclidean group is a Lie group is quoted here;
it is not needed for the matrix argument.
\end{beauty}

\section{Homomorphisms and the adjoint representation}

A group homomorphism must carry one-parameter subgroups to one-parameter
subgroups. Their initial velocities define its linear version.

\begin{theorem}\label{thm:homlie}
For a continuous homomorphism $\Phi:G\to H$ of matrix Lie groups there
is a unique real-linear map $\varphi:\mathfrak g\to\mathfrak h$ with
$\Phi(e^X)=e^{\varphi(X)}$ for every $X\in\mathfrak g$.
It satisfies
\[
 \varphi([X,Y])=[\varphi(X),\varphi(Y)],\qquad
 \varphi(gXg^{-1})=\Phi(g)\varphi(X)\Phi(g)^{-1}.
\]
The map $\Phi$ is smooth and $\varphi=d\Phi_I$.
\end{theorem}
\begin{proof}
For each $X$, Proposition~\ref{prop:oneparameter} gives a unique
$\varphi(X)$ with $\Phi(e^{tX})=e^{t\varphi(X)}$ for all $t$.
These exponentials belong to $H$, so $\varphi(X)\in\mathfrak h$.
Rescaling the parameter gives $\varphi(sX)=s\varphi(X)$ for every real $s$.
The product formula and continuity give
\[
 e^{t\varphi(X+Y)}
 =\Phi\left(\lim_m(e^{tX/m}e^{tY/m})^m\right)
 =e^{t(\varphi(X)+\varphi(Y))}.
\]
Differentiate at $t=0$ to obtain additivity.
The commutator formula similarly gives
$e^{t^2\varphi([X,Y])}=e^{t^2[\varphi(X),\varphi(Y)]}$.
Comparing second derivatives at zero yields the bracket identity.
Conjugating the one-parameter subgroup by $g$ and applying $\Phi$
gives the second identity in the same way, by first derivatives at zero.

A real-linear map between finite-dimensional spaces is smooth.
In the exponential charts of Theorem~\ref{thm:closedsubgroup},
$\Phi=\exp\circ\varphi\circ\log$ near $I$, so $\Phi$ is smooth there.
The identity
$\Phi=L_{\Phi(g)}\circ\Phi\circ L_{g^{-1}}$ near $g$
then proves smoothness everywhere. Differentiating along $e^{tX}$ gives
$d\Phi_I(X)=\varphi(X)$. Any linear map with the displayed exponential
property has the same derivatives after replacing $X$ by $tX$,
which proves uniqueness.
\end{proof}

\begin{proposition}\label{prop:connectedgen}
A connected matrix Lie group $G$ is generated by $\exp(\mathfrak g)$.
Every continuous homomorphism from $G$ is determined by its differential
at the identity.
\end{proposition}
\begin{proof}
The subgroup generated by $\exp(\mathfrak g)$ contains a neighbourhood
of $I$ by the local chart theorem. Translating that neighbourhood shows
this subgroup is open. Every other coset is open too, so its complement
is open. Connectedness forces it to be all of $G$.
The differential determines a homomorphism on every exponential by
Theorem~\ref{thm:homlie}, and hence on every product of exponentials.
\end{proof}

Conjugation gives a linear action on $\mathfrak g$:
$\mathrm{Ad}_g(X)=gXg^{-1}$.
It is called the adjoint representation. Its differential is
\[
 \mathrm{ad}_X(Y)=\left.\frac{\dd}{\dd t}e^{tX}Ye^{-tX}\right|_{t=0}
 =[X,Y].
\]
The homomorphism theorem therefore gives
$\mathrm{Ad}_{e^X}=e^{\mathrm{ad}_X}$.
These are linear maps on $\mathfrak g$, including the trivial action when
$\mathfrak g$ is zero-dimensional.

\section{Rotations and the spin group}

Every antisymmetric $3\times3$ matrix has the form
\[
 L(u)=\begin{pmatrix}0&-u_3&u_2\\u_3&0&-u_1\\-u_2&u_1&0\end{pmatrix},
 \qquad L(u)v=u\times v.
\]
The vector triple-product identity gives
$[L(u),L(v)]=L(u\times v)$. Thus the Lie algebra is $\R^3$ with its
cross-product bracket. If $u\ne0$, $L(u)$ vanishes on $\R u$ and acts
on $u^\perp$ as a quarter-turn times $|u|$.
Its exponential fixes the axis and rotates the perpendicular plane by
$|u|$. For $u=0$ the exponential is $I$.
Euler's rotation lemma (Volume~3, Lemma~12.3) says every element of
$SO(3)$ has an axis and angle, so every such element is $e^{L(u)}$.
The paths $t\mapsto e^{tL(u)}$ show that $SO(3)$ is path-connected.

For $SU(2)$, use the Pauli matrices
\[
 \sigma_1=\begin{pmatrix}0&1\\1&0\end{pmatrix},\quad
 \sigma_2=\begin{pmatrix}0&-i\\i&0\end{pmatrix},\quad
 \sigma_3=\begin{pmatrix}1&0\\0&-1\end{pmatrix}.
\]
Its Lie algebra consists of
\[
 X(u)=-\frac i2\sum_j u_j\sigma_j
 =\frac12\begin{pmatrix}-iu_3&-u_2-iu_1\\u_2-iu_1&iu_3\end{pmatrix}.
\]
Direct multiplication gives $\sigma_j^2=I$ and
$\sigma_1\sigma_2=i\sigma_3$, with cyclic versions and opposite signs
in the reverse orders. It follows that
\[
 [X(u),X(v)]=X(u\times v),\qquad X(u)^2=-|u|^2I/4.
\]
So $X(u)\mapsto L(u)$ is an isomorphism of Lie algebras.
The corresponding groups need not be isomorphic.
Indeed, the first column determines an $SU(2)$ matrix:
\[
 SU(2)=\left\{\begin{pmatrix}a&-\overline b\\b&\overline a\end{pmatrix}:
 |a|^2+|b|^2=1\right\}\cong S^3.
\]
The orthogonal second unit column has this form by the determinant-one
condition. These are also the unit quaternions of Volume~3.

\begin{theorem}[the spin double cover]\label{thm:spin}
In the basis $X(e_1),X(e_2),X(e_3)$, the adjoint representation of
$SU(2)$ is a continuous surjective homomorphism
$\pi:SU(2)\to SO(3)$ with kernel $\{I,-I\}$.
Its differential is the isomorphism $X(u)\mapsto L(u)$.
It is a two-sheeted universal covering. In particular,
$SO(3)\cong\R P^3$ as topological spaces and
$\pi_1(SO(3))\cong\Z/2$.
For each unit vector $u$, the full-turn loop
$\theta\mapsto e^{\theta L(u)}$, $0\le\theta\le2\pi$,
is not nullhomotopic, while the loop of two full turns is nullhomotopic.
\end{theorem}
\begin{proof}
\emph{The image consists of rotations.}
The real inner product $\langle X,Y\rangle=-2\operatorname{tr}(XY)$
makes the displayed basis orthonormal, since
$\operatorname{tr}(\sigma_j\sigma_k)=2\delta_{jk}$.
Conjugation preserves this inner product by cyclicity of trace.
Thus $\pi(g)$ is orthogonal. Its determinant is a continuous function
from connected $S^3$ to $\{1,-1\}$, equal to one at $I$;
it is therefore always one.

\emph{The differential and surjectivity.}
The differential is $\mathrm{ad}$, and
$\mathrm{ad}_{X(u)}X(v)=X(u\times v)$ represents $L(u)$.
Hence $\pi(e^{X(u)})=e^{L(u)}$.
Every spatial rotation has this exponential form, so $\pi$ is onto.

\emph{The kernel and covering.}
If $\pi(g)=I$, then $g$ commutes with all $X(u)$ and hence with all
Pauli matrices. Together with $I$ these span $M_2(\C)$, so $g$ is
scalar. Its determinant is one, forcing $g=I$ or $-I$.
Both do lie in the kernel.
The resulting map $S^3/\{\pm1\}\to SO(3)$ is a continuous bijection
from a compact space to a Hausdorff space, hence a homeomorphism.
The source is $\R P^3$, and $S^3\to\R P^3$ is the two-sheeted
covering of Example~\ref{ex:coverings}(d).
Since $S^3$ is simply connected by Corollary~\ref{cor:spheres}, this
is the universal covering. The covering-space calculation from
Example~\ref{ex:computations}(b) gives fundamental group $\Z/2$.

\emph{Lift a turn.}
Let $|u|=1$. Since $X(u)^2=-I/4$, separating even and odd powers gives
\[
 e^{\theta X(u)}=\cos(\theta/2)I+2\sin(\theta/2)X(u).
\]
This is the lift starting at $I$ of $e^{\theta L(u)}$.
At $\theta=2\pi$ it equals $-I$, so the lift is not closed.
A nullhomotopic loop would have a closed lift, by
Proposition~\ref{prop:fibre}; the full turn is therefore not nullhomotopic.
At $4\pi$ the lift returns to $I$. Its closed loop contracts in $S^3$,
and projection of that contraction contracts the two-turn loop in $SO(3)$.
\end{proof}

\begin{figure}[ht]
\centering
\fig{ch07-ballmodel}
\caption{$SO(3)$ as a ball of radius $\pi$: the point $\theta u$ is the rotation by $\theta$ about $u$, and antipodal points of the boundary sphere are glued (a rotation by $\pi$ about $u$ equals one about
$-u$). Left: a full turn about an axis is a diameter, a closed loop that cannot be shrunk; its two ends are glued, and moving one end drags the other to the opposite side. Right: two turns are two
diameters; rotating the second one until it meets the first, then pulling the pair inwards, shrinks the loop to a point. This is Dirac's belt trick.}
\label{fig:ballmodel}
\end{figure}

\begin{remark}
The general statement $\pi_1(SO(n))\cong\Z/2$ for $n\ge3$, and its
Clifford-algebra construction of the double cover $\mathrm{Spin}(n)$,
are quoted extensions. The proof here establishes the case $n=3$.
In two dimensions $SO(2)\cong S^1$ has fundamental group $\Z$,
as already proved for the circle. A two-turn planar loop still has
winding number two; the contraction above uses three-dimensional rotations.
\end{remark}

\section{$SL_2(\R)$ and the maximal torus of $U(n)$}

For $\mathfrak{sl}_2(\R)$ take the basis
\[
 H=\begin{pmatrix}1&0\\0&-1\end{pmatrix},\quad
 E=\begin{pmatrix}0&1\\0&0\end{pmatrix},\quad
 F=\begin{pmatrix}0&0\\1&0\end{pmatrix}.
\]
Multiplication gives $[H,E]=2E$, $[H,F]=-2F$ and $[E,F]=H$.
This algebra is not isomorphic to $\mathfrak{so}(3)$.
The map $\mathrm{ad}_H$ has the real eigenvalue two on $E$,
whereas every $\mathrm{ad}_{L(u)}$ acts as a cross product and has
only eigenvalues $0,\pm i|u|$.
Chapter~9 will interpret $e^{tH}$, $e^{tE}$ and $e^{t(E-F)}$ as
hyperbolic, parabolic and elliptic motions.
Exercise~\ref{ex:slexp} shows that not every $SL_2(\R)$ matrix is a
single exponential, even though products of exponentials can reach it.

For $U(n)$ the spectral theorem makes the picture particularly direct.
A diagonal unitary matrix has one unit-circle coordinate for each diagonal
entry. The group of these matrices is a torus, a product of $n$ circles.

\begin{theorem}[maximal torus and Weyl group of $U(n)$]\label{thm:torus}
Let $T\subset U(n)$ be the diagonal unitary subgroup, isomorphic to
$(S^1)^n$.
\begin{enumerate}
\item[(a)] Each element of $U(n)$ is conjugate by a unitary matrix to an
element of $T$. The exponential $\mathfrak u(n)\to U(n)$ is onto and
$U(n)$ is connected.
\item[(b)] If $t\in T$ has distinct diagonal entries, its centraliser in
$U(n)$ is exactly $T$. Consequently $T$ is maximal abelian.
\item[(c)] Its normaliser $N(T)=\{g\in U(n):gTg^{-1}=T\}$ consists
of the matrices $P_wD$, where $D\in T$ and $P_w$ is a permutation
matrix. The quotient $N(T)/T$ is $S_n$, acting by permuting diagonal entries.
\item[(d)] Two elements of $T$ are conjugate in $U(n)$ exactly when their
diagonal entries agree after a permutation.
\end{enumerate}
\end{theorem}
\begin{proof}
(a) A unitary matrix is normal. The spectral theorem
(Volume~8, Theorem~3.2) writes it as
$g=U\operatorname{diag}(e^{i\theta_j})U^*$.
Then $g=e^X$ with $X=U\operatorname{diag}(i\theta_j)U^*\in\mathfrak u(n)$,
and $s\mapsto e^{sX}$ joins $I$ to $g$.
(b) A matrix commuting with $t$ preserves each of its one-dimensional
coordinate eigenspaces. It is diagonal, and if unitary belongs to $T$.
Conversely all elements of $T$ commute. Any abelian subgroup containing
$T$ must commute with this particular $t$, so it cannot be larger.
(c) For $g\in N(T)$, the diagonal matrix $gtg^{-1}$ has the same
set of distinct eigenvalues as $t$. Some permutation $w$ gives
$gtg^{-1}=P_wtP_w^{-1}$. Thus $P_w^{-1}g$ commutes with $t$
and belongs to $T$, proving $g=P_wD$.
Conversely each such matrix normalises all diagonal matrices.
Its permutation of coordinate lines is unique, and composition composes
these permutations. This gives a homomorphism onto $S_n$ with kernel $T$.
(d) Conjugation preserves eigenvalues with multiplicity. Conversely
permutation matrices realise every permutation, including when entries repeat.
\end{proof}

\begin{figure}[ht]
\centering
\fig{v9-ch07-spin-lift}
\caption{A full-turn loop in $SO(3)$ lifts from $I$ to $-I$ in $SU(2)$; a second turn closes the lift. The upper and lower paths are schematic parameter paths, not distances in the groups. The endpoints have the same spatial rotation but different points of the covering sphere.}
\label{fig:v9-spin-lift}
\end{figure}
\begin{figure}[ht]
\centering
\fig{v9-ch07-affine-product}
\caption{The affine maps $T(x)=x+1$ and $D(x)=2x$ do not commute: $D(T(0))=2$ while $T(D(0))=1$. Their matrices lie in a two-dimensional Lie group. The order of the arrows records composition; the bracket of its dilation and translation directions is a translation direction.}
\label{fig:v9-affine-product}
\end{figure}

\section{What we proved, and what comes next}

The same tangent algebra can belong to different groups.
The map from $SU(2)$ to $SO(3)$ identifies their tangent spaces, but
identifies two points of the sphere with every spatial rotation.
Following a loop detects that identification, which a derivative at one
point cannot see.

A closed matrix subgroup has smooth charts, its one-parameter motions
are exponentials, and continuous homomorphisms respect their brackets.
Connectedness then lets a differential determine a homomorphism on the
whole group. It does not make the exponential itself injective or onto
in every example; the exercises let you test both issues explicitly.

For $U(n)$, conjugacy leaves only the unordered eigenvalues on the circle.
Chapter~8 asks for their distribution under a uniform random choice of
matrix. Haar measure supplies that choice, and Weyl's formula describes
how it changes when we pass to eigenvalue coordinates.

\section{Exercises}

\begin{exercise}[$\star\star$]\label{ex:slexp}
\emph{Which matrices are single exponentials?}
(a) For $X\in\mathfrak{sl}_2(\R)$, prove $X^2=-\det(X)I$ and
compute $e^X$ for negative, zero and positive $\det X$.
(b) Prove $\operatorname{tr}e^X\ge-2$, with equality only if $e^X=-I$.
(c) Show that $\left(\begin{smallmatrix}-2&0\\0&-1/2\end{smallmatrix}\right)$
and $\left(\begin{smallmatrix}-1&1\\0&-1\end{smallmatrix}\right)$
are not single exponentials of traceless real matrices but are products
of two such exponentials.
\end{exercise}
\begin{exercise}[$\star$]\label{ex:heisenberg}
\emph{A group with a finite exponential series.}
For $G=\{\left(\begin{smallmatrix}1&a&c\\0&1&b\\0&0&1\end{smallmatrix}\right):a,b,c\in\R\}$,
(a) find its Lie algebra and brackets, and show that its commutator
algebra is central; (b) prove the exponential is bijective and
$e^Xe^Y=e^{X+Y+[X,Y]/2}$;
(c) explain the algebraic relation to the position--momentum commutator
$[Q,P]=i\hbar I$. For a concrete calculation you may use the operators
$Qf=xf$ and $Pf=-i\hbar f'$ on complex polynomials.
\end{exercise}
\begin{exercise}[$\star$]\label{ex:oneparamso3}
\emph{Axes and constant angular velocity.}
Describe all continuous homomorphisms $\R\to SO(3)$ as rotations about
a fixed unit axis at constant angular speed. Determine their periods,
including the constant case.
Show that two nonidentity rotations commute exactly when they have the
same axis or are half-turns about perpendicular axes. Explain the identity
exception explicitly.
\end{exercise}
\begin{exercise}[$\star\star$]\label{ex:sondconnected}
\emph{Orthogonal groups in dimension $n\ge1$.}
(a) Put each $A\in SO(n)$ into block-diagonal form by conjugation in
$SO(n)$, using plane rotations and scalar blocks $1,-1$, with an even
number of $-1$ blocks. Deduce exponential surjectivity and connectedness.
(b) Prove $O(n)$ has exactly two connected components.
What changes if zero-size matrices are allowed and $n=0$?
\end{exercise}
\begin{exercise}[$\star\star$]\label{ex:sl2pi1}
\emph{Rotations and stretches in $SL_2(\R)$.}
Use polar decomposition (Volume~8, Theorem~4.3) to prove
\[
 SL_2(\R)\cong SO(2)\times\{P=P^{\mathsf T}>0:\det P=1\}
 \cong S^1\times\R^2.
\]
Deduce $\pi_1(SL_2(\R))\cong\Z$.
A quoted advanced fact, unused in this exercise, is that its universal
cover has no faithful finite-dimensional matrix representation.
\end{exercise}
\begin{exercise}[$\star$]\label{ex:centersu2}
\emph{Centres distinguish groups.}
Find the centres of $SU(2)$ and $SO(3)$ and use them to prove that the
groups are not isomorphic, despite their isomorphic Lie algebras.
\end{exercise}

\begin{exercise}[$\star\star$]\label{ex:v9-rodrigues}
\emph{A rotation formula and its ambiguity.}
For $|u|=1$, put $A=L(u)$. Show that $A^2=uu^{\mathsf T}-I$ and
\[
 e^{\theta A}=I+\sin\theta\,A+(1-\cos\theta)A^2.
\]
Deduce $\operatorname{tr}e^{\theta A}=1+2\cos\theta$.
Explain how to recover the angle in $[0,\pi]$ and its axis from a rotation
matrix, treating angle zero and angle $\pi$ separately.
\end{exercise}
\begin{exercise}[$\star\star$]\label{ex:v9-circle-homomorphisms}
\emph{A differential with an integer constraint.}
Classify all continuous homomorphisms $\chi:U(1)\to U(1)$.
Which real-linear maps $i\R\to i\R$ occur as their differentials?
Then prove that every continuous homomorphism $SU(2)\to U(1)$ is trivial.
\end{exercise}
\begin{exercise}[$\star\star$]\label{ex:v9-affine-exponential}
\emph{Translations and dilations.}
Let $G=\{\left(\begin{smallmatrix}a&b\\0&1\end{smallmatrix}\right):a>0,b\in\R\}$.
Show it is closed in $GL_2(\R)$. Find its Lie algebra and bracket.
Compute the exponential and prove it is a bijection from the Lie algebra
onto $G$, giving a formula for its inverse. Interpret the two coordinate
motions as maps of a real line.
\end{exercise}
\begin{exercise}[$\star\star$]\label{ex:v9-spin-traces}
\emph{A half-turn upstairs and downstairs.}
For $g\in SU(2)$, prove
$\operatorname{tr}\pi(g)=(\operatorname{tr}g)^2-1$.
Describe the preimage of the half-turns as a subset of $S^3$.
For a unit vector $u$ and integer $k$, determine exactly when the loop
$\theta\mapsto e^{k\theta L(u)}$, $0\le\theta\le2\pi$, contracts.
\end{exercise}

\section{Hints}

\begin{hint}{ex:slexp}
Use Cayley--Hamilton, separate even and odd powers, and multiply by
$-I=e^{\pi(E-F)}$ for the final constructions.
\end{hint}
\begin{hint}{ex:heisenberg}
Use $E_{12},E_{23},E_{13}$ as a basis. All products of three strictly
upper triangular $3\times3$ matrices vanish.
\end{hint}
\begin{hint}{ex:oneparamso3}
A commuting rotation must preserve the fixed axis as a line. It may
reverse its unit direction; consider that case separately.
\end{hint}
\begin{hint}{ex:sondconnected}
Pair nonreal conjugate eigenvalues into real rotation planes and pair
$-1$ eigenvalues into half-turn planes.
\end{hint}
\begin{hint}{ex:sl2pi1}
For $M=A^{\mathsf T}A$ of determinant one, verify
$P=(M+I)/\sqrt{\operatorname{tr}M+2}$.
A positive-definite determinant-one symmetric matrix is determined by its
upper-left entry $a>0$ and off-diagonal entry $b$.
\end{hint}
\begin{hint}{ex:centersu2}
A central element commutes with every one-parameter subgroup, hence with
every tangent matrix. Use Pauli matrices in one case and rotation axes
in the other.
\end{hint}
\begin{hint}{ex:v9-rodrigues}
The triple-product identity gives $A^3=-A$. Group the exponential series
into its odd and even terms after the identity term.
\end{hint}
\begin{hint}{ex:v9-circle-homomorphisms}
Apply the one-parameter theorem to $t\mapsto\chi(e^{it})$.
For $SU(2)$, every basis vector of its Lie algebra is a bracket.
\end{hint}
\begin{hint}{ex:v9-affine-exponential}
For $X=\left(\begin{smallmatrix}s&t\\0&0\end{smallmatrix}\right)$,
compute $X^k$ and extend $(e^s-1)/s$ continuously at $s=0$.
\end{hint}
\begin{hint}{ex:v9-spin-traces}
Write $g=rI+N$ with $r=\operatorname{tr}g/2\in[-1,1]$.
Cayley--Hamilton gives $N^2=-(1-r^2)I$. Handle $r=\pm1$ separately.
\end{hint}

\section{Solutions}

\begin{solution}{ex:slexp}
\idea A traceless two-by-two matrix squares to a scalar, leaving only two
series to sum.
\writeup (a) Cayley--Hamilton gives $X^2=-\det(X)I$.
For $\det X=-\lambda^2<0$, the even and odd powers sum to
$e^X=\cosh\lambda\,I+(\sinh\lambda/\lambda)X$.
For $\det X=\mu^2>0$, they give
$e^X=\cos\mu\,I+(\sin\mu/\mu)X$.
For $\det X=0$, all powers from $X^2$ onward vanish and $e^X=I+X$.
(b) Their traces are $2\cosh\lambda$, $2\cos\mu$ and $2$.
They are at least $-2$. Equality requires $\cos\mu=-1$, giving
$\sin\mu=0$ and $e^X=-I$.
(c) The negative diagonal matrix has trace $-5/2$, and the second
matrix has trace $-2$ without being $-I$, so neither is a single
exponential. Direct multiplication gives the products
\[
 \begin{pmatrix}-2&0\\0&-1/2\end{pmatrix}
 =e^{\pi(E-F)}e^{(\log2)H},\qquad
 \begin{pmatrix}-1&1\\0&-1\end{pmatrix}
 =e^{\pi(E-F)}e^{-E}.
\]
All exponents are traceless real matrices, as required.
\end{solution}
\begin{solution}{ex:heisenberg}
\idea Nilpotence turns the exponential and product identities into finite
polynomial calculations.
\writeup (a) Differentiating the defining entries shows that the algebra
is spanned by $A=E_{12}$, $B=E_{23}$, $Z=E_{13}$.
Conversely their strictly upper triangular combinations exponentiate
into $G$. The only nonzero basic bracket is $[A,B]=Z$;
$Z$ commutes with all three, so the commutator algebra is central.
(b) For $X=xA+yB+zZ$, $X^2=xyZ$ and $X^3=0$.
Thus $e^X=I+xA+yB+(z+xy/2)Z$.
Given entries $a,b,c$, the unique preimage has
$x=a,y=b,z=c-ab/2$.
For two such matrices, every product of three vanishes, so
\[
 e^Xe^Y=I+X+Y+X^2/2+XY+Y^2/2.
\]
Expanding $e^{X+Y+[X,Y]/2}$ gives the same expression because
$[X,Y]$ is a multiple of $Z$ and its products with $X,Y$ vanish.
(c) On complex polynomials,
$[Q,P]f=-i\hbar xf'+i\hbar(xf)'=i\hbar f$.
The assignments $A\mapsto Q$, $B\mapsto P$, $Z\mapsto i\hbar I$
therefore respect exactly the brackets just computed.
This is an algebraic operator realisation on an infinite-dimensional
space. For $\hbar\ne0$ it cannot hold for finite-dimensional complex
matrices: the trace of a commutator is zero, while
$\operatorname{tr}(i\hbar I)=i\hbar d\ne0$ in dimension $d\ge1$.
The quantum-mechanical relation motivates the name, without making the
three-by-three matrices themselves position and momentum operators.
\end{solution}
\begin{solution}{ex:oneparamso3}
\idea The generator fixes the angular velocity; commutation preserves its
axis but may reverse its direction.
\writeup Every homomorphism is $e^{tL(v)}$.
For $v\ne0$, set $u=v/|v|$ and $\omega=|v|$; it is rotation about $u$
through angle $\omega t$, with least positive period $2\pi/\omega$.
For $v=0$ it is constant and every positive time is a period, so there
is no least positive period.

Let nonidentity $R,S$ commute, and let $u$ be a unit axis of $R$.
Its fixed space is exactly $\R u$, including when $R$ is a half-turn.
Thus $Su=u$ or $Su=-u$.
In the first case, $S$ has the same axis. In the second, $-1$ is an
eigenvalue of $S$, so it is a half-turn about an axis perpendicular to
$u$. Conjugation by $S$ changes a rotation of angle $\theta$ about $u$
into one of angle $-\theta$ about $u$.
Equality with nonidentity $R$ therefore forces $\theta=\pi$ modulo $2\pi$.
Conversely rotations about one axis commute. Two perpendicular half-turns
are diagonal matrices with entries $1,-1,-1$ in a suitable orthonormal
basis, so they commute as well.
The identity commutes with every rotation, without any axis restriction.
\end{solution}
\begin{solution}{ex:sondconnected}
\idea Real orthogonal matrices split into invariant lines and rotation planes.
\writeup (a) The complex spectral theorem gives orthogonal eigenspaces
of a unitary matrix. Nonreal eigenvalues of a real orthogonal matrix
come in conjugate pairs $e^{i\theta},e^{-i\theta}$.
If $v=x+iy$ is an eigenvector, orthogonality to $\overline v$ gives
$|x|=|y|$ and $x\cdot y=0$. Their real span is an invariant plane with
an orthonormal basis in which $A$ is a plane rotation.
The real orthogonal complement of each such plane is invariant too,
since $A$ is orthogonal and preserves the plane. Repeat on that complement;
this includes repeated eigenvalues. The remaining real eigenvalues are
$1$ or $-1$. Together these give a real orthonormal block basis.
The determinant condition makes the number of $-1$ blocks even.
Pair them into plane rotations through $\pi$.
If the basis change has determinant $-1$, reverse one basis vector;
this only reverses a rotation angle or leaves a scalar block unchanged,
and makes the change belong to $SO(n)$.
Each rotation block is the exponential of a skew-symmetric two-by-two
block; scalar $1$ blocks use zero. Conjugating their direct sum back
gives $A=e^X$, $X\in\mathfrak{so}(n)$.
The path $e^{tX}$ proves connectedness, including the trivial case $n=1$.
(b) For $n\ge1$ both determinant signs occur, and the continuous
sign map separates them. The positive fibre is connected by (a).
Multiplying by a fixed reflection identifies it with the negative fibre,
which is therefore also connected. There are exactly two components.
For $n=0$, both groups consist only of the empty identity matrix,
whose determinant is one; there is only one component.
\end{solution}
\begin{solution}{ex:sl2pi1}
\idea Give explicit continuous coordinates for the positive stretch, so
no choice of eigenvectors or contour formula is needed.
\writeup Let $M=A^{\mathsf T}A$. It is positive definite with determinant
one. Cayley--Hamilton gives $M^2-(\operatorname{tr}M)M+I=0$, hence
\[
 P=\frac{M+I}{\sqrt{\operatorname{tr}M+2}},\qquad P^2=M.
\]
This $P$ is symmetric positive definite, has determinant one, and depends
continuously on $A$. Then $W=AP^{-1}$ is orthogonal with determinant
one. These are the unique polar factors, and the inverse operation is
the continuous product $(W,P)\mapsto WP$.

Every symmetric positive-definite determinant-one matrix is uniquely
\[
 P(a,b)=\begin{pmatrix}a&b\\b&(1+b^2)/a\end{pmatrix},
 \qquad a>0,\quad b\in\R.
\]
The leading entry is positive. Conversely the associated form is
$a(x+by/a)^2+y^2/a$, which is positive on every nonzero pair.
Thus all these parameters give positive-definite matrices.
Writing $a=e^x$, $b=y$ gives continuous coordinates $(x,y)\in\R^2$
with continuous inverse $x=\log P_{11}$, $y=P_{12}$.
Thus the polar decomposition is the claimed homeomorphism.
Contract $(x,y)$ linearly to $(0,0)$ while keeping $W$ fixed.
This deformation retracts onto $SO(2)\cong S^1$ and gives fundamental
group $\Z$, generated by one full circle of plane rotations.
The quoted non-matrix property of the universal cover is not used.
\end{solution}
\begin{solution}{ex:centersu2}
\idea A central matrix must commute with every infinitesimal rotation.
\writeup If $g$ is central, differentiate $ge^{tX}=e^{tX}g$ at zero.
For $SU(2)$ this makes $g$ commute with all Pauli matrices, hence with
all complex matrices, so it is scalar. Determinant one gives $g=\pm I$,
and both are central.
For $SO(3)$, the differentiated condition says
$gL(u)g^{-1}=L(u)$ for every $u$.
Rotations preserve the cross product, so the left side is $L(gu)$.
Injectivity of $L$ gives $gu=u$ for every $u$, hence $g=I$.
An isomorphism preserves the centre, but these centres have different
sizes. The two groups are therefore not isomorphic as groups.
\end{solution}

\begin{solution}{ex:v9-rodrigues}
\idea Powers of the generator reduce to $I,A,A^2$.
\writeup The vector triple product gives
$A^2v=u(u\cdot v)-v$, so $A^2=uu^{\mathsf T}-I$ and $A^3=-A$.
Odd powers after $A$ alternate between $A$ and $-A$;
even powers after $I$ alternate between $A^2$ and $-A^2$.
Summing their sine and cosine series gives the required formula.
Since $\operatorname{tr}A=0$ and $\operatorname{tr}A^2=-2$,
the trace is $1+2\cos\theta$.
For any rotation $R$, choose its angle $\theta\in[0,\pi]$ by
$\cos\theta=(\operatorname{tr}R-1)/2$.
If $0<\theta<\pi$, then
$L(u)=(R-R^{\mathsf T})/(2\sin\theta)$ recovers the oriented unit axis.
If $\theta=0$, $R=I$ and every axis works.
If $\theta=\pi$, use the one-dimensional $+1$ eigenspace of $R$;
its two unit directions $u,-u$ give the same half-turn.
No division by $\sin\theta$ is made at these endpoints.
\end{solution}
\begin{solution}{ex:v9-circle-homomorphisms}
\idea Periodicity restricts an otherwise arbitrary infinitesimal angular speed.
\writeup The composite $t\mapsto\chi(e^{it})$ is a continuous
one-parameter subgroup of $U(1)$, so equals $e^{ict}$ for some real $c$.
Since the input has period $2\pi$, $e^{2\pi ic}=1$, forcing $c=k\in\Z$.
Thus $\chi(z)=z^k$ for all $z\in U(1)$; each integer power is indeed a
continuous homomorphism. Their differentials are precisely $it\mapsto ikt$.
A general real-linear map would allow any real $c$, but only integers
respect the closed circle.

For $\Phi:SU(2)\to U(1)$, its differential annihilates brackets because
the target algebra $i\R$ is abelian. The relations
$[X(e_1),X(e_2)]=X(e_3)$ and their cyclic versions show that brackets
span the whole source algebra. Thus $d\Phi_I=0$.
Connectedness of $SU(2)$ and Proposition~\ref{prop:connectedgen} force
$\Phi$ to be the trivial homomorphism.
\end{solution}
\begin{solution}{ex:v9-affine-exponential}
\idea The dilation parameter exponentiates; the translation parameter is
scaled by the average of that dilation.
\writeup Products and inverses preserve the displayed form with $a>0$.
A limit in $GL_2(\R)$ still has bottom row $(0,1)$ and nonzero top-left
entry. A limit of positive such entries must remain positive, so $G$
is relatively closed. Differentiating gives
$X(s,t)=\left(\begin{smallmatrix}s&t\\0&0\end{smallmatrix}\right)$,
and conversely its exponential lies in $G$.
Multiplication yields
$[X(s,t),X(s',t')]=X(0,st'-s't)$.
For $k\ge1$, the top row of $X^k$ is $(s^k,s^{k-1}t)$.
Summing gives
\[
 e^{X(s,t)}=\begin{pmatrix}e^s&t\,\eta(s)\\0&1\end{pmatrix},
 \qquad \eta(s)=\begin{cases}(e^s-1)/s,&s\ne0,\\1,&s=0.\end{cases}
\]
The series shows $\eta$ is continuous and smooth at zero, and it is
positive for every real $s$.
For a given $(a,b)$, the unique inverse parameters are
$s=\log a$, $t=b/\eta(\log a)$, proving bijectivity.
Acting on $(x,1)^{\mathsf T}$ gives $x\mapsto ax+b$.
The generators $X(1,0)$ and $X(0,1)$ give dilations $x\mapsto e^s x$
and translations $x\mapsto x+t$; their bracket is the translation
generator, consistent with Figure~\ref{fig:v9-affine-product}.
\end{solution}
\begin{solution}{ex:v9-spin-traces}
\idea The half-angle upstairs becomes a full angle under the covering map.
\writeup The $SU(2)$ form gives real trace $2r$ with $r\in[-1,1]$.
If $r=\pm1$, the unit-sphere coordinates force $g=\pm I$;
then $\pi(g)=I$ and the trace identity reads $3=4-1$.
Otherwise put $N=g-rI$.
Cayley--Hamilton gives $N^2=-(1-r^2)I$.
Also $g^{-1}=2rI-g=g^*$, so $N^*=-N$ and $\operatorname{tr}N=0$.
Write $N=2\sqrt{1-r^2}X(u)$; the square identity for $X(u)$ forces
$|u|=1$. Choose $t\in(0,\pi)$ with $\cos t=r$.
Then $g=e^{2tX(u)}$ and $\pi(g)=e^{2tL(u)}$.
Its trace is $1+2\cos(2t)=4\cos^2t-1=(\operatorname{tr}g)^2-1$.
A spatial half-turn has trace $-1$, so its preimage is exactly
$\operatorname{tr}g=0$. In the coordinates $(a,b)\in S^3\subset\C^2$,
this is $\operatorname{Re}a=0$, a great two-sphere.

The $k$-turn loop lifts from $I$ to
$e^{2\pi kX(u)}=(-1)^kI$.
For odd $k$ it cannot contract because the lift is not closed.
For even $k$ its lift is a closed loop in simply connected $S^3$,
so a contraction projects to a contraction downstairs.
This includes $k=0$ and negative integers: precisely the even integers
represent the trivial loop class.
\end{solution}

\section*{Sources and further reading}
\addcontentsline{toc}{section}{Sources and further reading}

G.~E.~Uhlenbeck and S.~Goudsmit, ``Ersetzung der Hypothese vom unmechanischen Zwang durch eine Forderung bez\"uglich des inneren Verhaltens jedes einzelnen Elektrons'', \emph{Naturwissenschaften}
13 (1925), 953--954; G.~E.~Uhlenbeck, ``Personal reminiscences'', \emph{Physics Today} 29, no.~6 (1976), 43--48, which tells the story of Lorentz's calculation and Ehrenfest's answer; W.~Pauli, ``Zur Quantenmechanik
des magnetischen Elektrons'', \emph{Z.\ Phys.} 43 (1927), 601--623. The neutron experiments: H.~Rauch et al., ``Verification of coherent spinor rotation of fermions'', \emph{Phys.\ Lett.\ A} 54 (1975), 425--427, and
S.~A.~Werner et al., ``Observation of the phase shift of a neutron due to precession in a magnetic field'', \emph{Phys.\ Rev.\ Lett.} 35 (1975), 1053--1055.

J.~von Neumann, ``\"Uber die analytischen Eigenschaften von Gruppen linearer Transformationen und ihrer Darstellungen'', \emph{Math.\ Z.} 30 (1929), 3--42; \'E.~Cartan, \emph{La th\'eorie des groupes finis et
continus et l'Analysis situs} (M\'emorial des sciences math\'ematiques 42, Gauthier-Villars, 1930). A.~Gleason, ``Groups without small subgroups'', \emph{Ann.\ of Math.} 56 (1952), 193--212, and D.~Montgomery and
L.~Zippin, ``Small subgroups of finite-dimensional groups'', same volume, 213--241.

The treatment follows B.~Hall, \emph{Lie Groups, Lie Algebras, and Representations}, 2nd ed. (Springer, 2015), Chapters~1--3 and~11, and W.~Rossmann, \emph{Lie Groups: An Introduction through Linear Groups}
(Oxford University Press, 2002). For the belt trick and its physics, R.~Penrose and W.~Rindler, \emph{Spinors and Space-Time}, vol.~1 (Cambridge University Press, 1984), Chapter~1.
